% Generated draft: requires mathematical and source review.
\documentclass[11pt]{article}
\usepackage[margin=25mm]{geometry}
\usepackage{fontspec}
\setmainfont{Latin Modern Roman}
\usepackage{amsmath,amssymb,mathtools,unicode-math}
\setmathfont{Latin Modern Math}
\usepackage{enumitem,needspace,xurl,hyperref}
\setlength{\parindent}{0pt}
\setlength{\parskip}{5pt}
\setlength{\emergencystretch}{3em}
\newcommand{\N}{\mathbb{N}}
\newcommand{\Z}{\mathbb{Z}}
\newcommand{\Q}{\mathbb{Q}}
\newcommand{\R}{\mathbb{R}}
\newcommand{\C}{\mathbb{C}}
\newcommand{\F}{\mathbb{F}}
\newcommand{\A}{\mathbb{A}}
\newcommand{\PP}{\mathbb{P}}
\newcommand{\E}{\mathbb{E}}
\newcommand{\eps}{\varepsilon}
\newcommand{\dd}{\,\mathrm{d}}
\begin{document}
\section*{UnsolvedMath ID: 8800035}
\textbf{Source problem code:} catalogue entry
\subsection{Definitions and mathematical statement}
For an elliptic curve \(E/\mathbb F_q\), let
\[
 e_q(E)=\min\{m\ge1:mP=0\text{ for every }P\in E(\mathbb F_q)\}
\]
be the exponent of its finite abelian group of rational points. Writing \(E(\mathbb F_q)\cong\mathbb Z/d_1\mathbb Z\times\mathbb Z/d_2\mathbb Z\), with \(d_1\mid d_2\), gives \(e_q(E)=d_2\). The source asks whether this exponent is typically close to \(q\).

The sampling and comparison are essential. One regime fixes the finite field and samples curves, then studies the resulting distributions as \(q\) grows; another fixes an elliptic curve over \(\mathbb Q\) and varies primes of good reduction. Sampling Weierstrass models, isomorphism classes or automorphism-weighted classes gives different statements. ``Close'' may mean ratio tending to \(1\), a lower bound \(q^{1-o(1)}\), or a weaker comparison. No one of these is specified by the excerpt.

The summary attributes an at-least-75-percent cyclicity result in a varying-curve regime to Vladut (1999). On cyclic instances \(e_q(E)=\#E(\mathbb F_q)\), whose ratio to \(q\) tends to \(1\) by the Hasse bound. It separately attributes to Duke (2003) the estimate \(e_p(E)\ge p^{3/4-\varepsilon}\) for almost all primes in a fixed-curve regime, and mentions ERH-conditional cyclicity work of Cojocaru--Murty--Duke. Exact sampling, curve hypotheses, constants and conditional statements are not supplied.

These reports do not specify one universal typical law. In particular a lower bound of exponent \(3/4-\varepsilon\) is not an upper bound proving that typical exponents are below \(p\), despite such language in the summary. Its ``often yes but not always'' conclusion is therefore retained only as a qualitative source assessment, not a resolution of an undefined probability/asymptotic question. The partial-progress status requires current primary-literature review.
\subsection{Short English statement}
Is the exponent of a typical elliptic-curve group nearly the field size? The answer must specify how curves or primes are sampled and how close is measured; the source reports results in distinct regimes without fixing a single asymptotic law.
\subsection{Sources}
\begin{itemize}
\item Local registry: \url{data/unsolvedmath/problems.json}, permanent UnsolvedMath ID \(8800035\), source code catalogue entry; supplied statement, background, and research summary, updated 2026-08-06. Attributed to Igor Shparlinski.
\item IPAM, \emph{Workshop on Number Theory and Cryptography---Open Problems} (2006), attachment 6224, PDF page 13: \url{https://www.ipam.ucla.edu/programs/workshops/workshop-i-number-theory-and-cryptography-open-problems/?tab=schedule}. Supplied PDF/legacy-PPT visual audit accessed 2026-07-29; it reports 19 inspected surviving attachments.
\item Vladut (1999); Duke (2003); Cojocaru--Murty--Duke conditional cyclicity work, as cited by the supplied summary. No full titles or exact theorem hypotheses are included.
\end{itemize}
\end{document}
